\documentclass[10pt]{article}
\usepackage[margin=0.8in]{geometry}
\usepackage{amsmath,amssymb}
\usepackage[hidelinks]{hyperref}
\usepackage{xurl}
\newcommand{\Tr}{\operatorname{Tr}}
\begin{document}
\begin{center}\textbf{The clock row: flux identity and selecting principle}\end{center}

\paragraph{Finite law and the question.}
Use one finite sine parent, its native CAR, all crossing matrix
elements and the declared positive ordinary-adjoint completion.
Rulers are prescribed commuting scalars, $0<s_-\leq\psi_x\leq s_+$;
all packet masses are positive. No continuum expansion is used.
With $A_s=m_s\Psi^{-1}+i\Psi^{-3}D_a$, define
\begin{align}
 E_{s0}&=|A_s|,\qquad E_{s1}=|A_s^\dagger|
                    =U_sE_{s0}U_s^\dagger,\notag\\
 t_x&=\operatorname{tr}_2\sum_s\eta_s(E_{s1}-E_{s0})_{xx},
       &\sum_xt_x&=0,\label{eq:fg-trace}\\
 H_*[N,\psi]&=\sum_xN_xC_x,\qquad
 C_x=\psi_x(t_x\mathbf1+\rho_x),
       &\rho_x&=\sum_rQ_{rx}^\dagger Q_{rx}\succeq0.
 \label{eq:fg-law}
\end{align}
Here $Q_{rx}$ includes the reference, moving filter, unit-lapse
denominator and complete dressing, all independent of $N$.
The formulas hold as finite matrices or positive forms with finite
expectations on their common domain. The trace equality follows
packet by packet from the common singular values; a resolvent
similarity gives the same equality. It does not assert $C_x=0$.
% [X]: polar singular values and resolvent similarity, at the stated
% finite scalar domain; see the accepted bridge proof, Sec. 1.

\paragraph{A Gauss form follows; a Gauss constraint does not.}
There are two explicit flux representations. For any connected
oriented graph on these sites, let $B$ be its outgoing incidence
matrix, $M$ a positive diagonal edge matrix and ${}^+$ the
Moore--Penrose inverse. Then
\[
 F=M B^T(BMB^T)^+(-t),\qquad BF=-t.
\]
The zero-sum identity in \eqref{eq:fg-trace} proves the second
equality. Edge cycles can be added freely; this construction is
generally nonlocal and depends on $M$.

A second route uses the same polar law directly:
$E_{s1}-E_{s0}=[U_s,E_{s0}U_s^\dagger]$. Put $V_s=E_{s0}U_s^\dagger$
and
\[
 j_{xy}=\sum_s\eta_s\Re\operatorname{tr}_2
       \bigl[(U_s)_{xy}(V_s)_{yx}-(V_s)_{xy}(U_s)_{yx}\bigr].
\]
Finite spin-trace cyclicity gives $j_{yx}=-j_{xy}$ and
$\sum_yj_{xy}=t_x$. Thus $F_{xy}=-j_{xy}$ is an exact, generally
dense crossing flux. In particular, for $L_x=N_x\psi_x$,
\[
 W_0=\sum_xL_xt_x=\frac12\sum_{x,y}(L_x-L_y)j_{xy}.
\]
Both routes express the vacuum density as a divergence without
imposing any equation on a state. The additional clock equation
$C_x=0$ would give
\begin{equation}
 (\operatorname{div}F)_x=\rho_x=e_x/\psi_x,
 \qquad e_x=\psi_x\rho_x.
 \label{eq:fg-gauss}
\end{equation}
For mean constraints the same statement uses expectations. Neither
representation identifies a unique local physical flux or a
canonical flux--readout commutator.
% [X]: graph pseudoinverse and native polar commutator routes.

For an open parent $\Lambda=I\sqcup B_\partial$ with
$\psi_b=1$, $N_b=N_\infty>0$, let
$H_\partial=\partial_{N_\infty}H_*=\sum_bC_b$, so
$H_*=N_\infty H_\partial+\sum_{i\in I}N_iC_i$.
The exact weighted trace yields
\[
 H_\partial=\sum_x\rho_x-\sum_{i\in I}C_i/\psi_i.
\]
Hence the boundary expectation is nonnegative on the mean clock
surface, retaining collar source. At unit external clock the
comparison lapse $N_x^*=\psi_x^{-1}$ independently gives the same
bound. On a closed parent imposing every mean row instead forces
$\sum_x\langle\rho_x\rangle=0$. These are conditional consequences
of the row, not arguments that every quantum state obeys it.
% [X]: weighted trace and constant-allocation comparison lapse.

\paragraph{A native counterfamily to weaker static premises.}
Take a periodic sine lattice with $V$ sites, uniform
$\psi=P>0$, $N=N_0>0$, and one massive light excitation at $p=0$.
Its protected zero/one-excitation law is free, and
\[
 E(p)=N_0\sqrt{m^2+P^{-4}a^{-2}\sum_i\sin^2(ap_i)},
 \qquad E(0)=N_0m.
\]
The sea energy is zero on every uniform background. Translation
covariance makes each local ruler expectation derivative equal;
their sum is $\partial_PE(0)=0$. The eigenstate therefore satisfies
state stationarity and $\langle\partial_{\psi_x}H_*\rangle=0$,
but
\begin{equation}
 \langle C_x\rangle=m/V>0.\label{eq:fg-counter}
\end{equation}
One check is $\sum_x\langle C_x\rangle=\partial_{N_0}E=m$
and translation invariance. Directly, the positive one-particle
block is $E_0^{1/2}NPE_0^{1/2}$; its local clock derivative is
$PE_0^{1/2}|x\rangle\langle x|E_0^{1/2}$ and
$E_0|p=0\rangle=(m/P)|p=0\rangle$, again giving $m/V$.
The vacuum clock derivative vanishes at this uniform background.
This is an unconstrained diagnostic, not a loaded periodic H2
sector. It proves that positive norms, positive excitation and
clock forms, lapse affinity, vacuum protection, state stationarity
and mean ruler stationarity do not imply the clock Gauss equation.
% [X]: dispersion/Hellmann--Feynman plus translation; direct positive
% block differentiation. No numerical control is used.

\paragraph{Two precise additional selecting principles.}
Set the external clock to $N_\infty=1$ in this paragraph.
If every compactly supported interior $N_x(t)$ is freely varied
as an auxiliary variable in the state action
\[
 I=\int dt\,\bigl(i\langle\Psi|\dot\Psi\rangle
       -\langle H_\partial+\sum_{x\in I}N_xC_x\rangle\bigr),
\]
with the exterior clock fixed, lapse variation gives
$\langle C_x\rangle=0$. Small variations of either sign are
allowed about a strictly positive lapse. This mean equation
requires an auxiliary-clock variational principle; self-adjoint
quantum evolution alone does not provide it. It does not imply
$C_x|\Psi\rangle=0$. Time reparametrization alone is weaker still:
the parametrized action
$\int(p\dot q+p_T\dot T-N[p_T+C])dt$ imposes $p_T+C=0$, permitting
$C\ne0$ unless a clock-momentum sector is separately fixed.

A strong constraint follows from a specified operational
redundancy. Assume a common invariant physical sector $\mathcal P$
for the self-adjoint generators at every interior lapse, on a
finite regulator or a common invariant domain. Require its
evolution operator, with fixed exterior clock, to equal boundary
evolution for every compact interior lapse change. Differentiating
a short time interval, or using the Duhamel insertion
\[
 \frac{\delta U(t_f,t_i)}{\delta N_x(t)}
       =-iU(t_f,t)C_xU(t,t_i),
\]
gives $C_x\mathcal P=0$. Conversely that invariant kernel removes
the interior lapse dependence. Equality only of physical rays
allows $C_x|_{\mathcal P}=c_x\mathbf1$; setting such sector-dependent
central charges to zero needs a normalization principle. With
$P_{\mathcal P}$ the sector's orthogonal projector, a compressed
equality $P_{\mathcal P}C_xP_{\mathcal P}=0$ without invariance
does not exclude leakage out of the sector.
% [X]: short-time generator differentiation and Duhamel insertion.

Thus the scalar law supplies an exact Gauss \emph{form}; a mean
constraint can be selected by auxiliary-clock stationarity, and a
strong constraint by the stated interior-clock redundancy with
zero central charge. These assumptions are weaker in scope than
the full H2 physical bridge, but are not consequences of the
counterexample's static premises. Identifying the physical sector,
its preservation and a native local flux algebra remains a
separate choice or theorem. No untraced operator partial trace or
noncommuting-ruler constraint is inferred from the scalar identity.

\paragraph{SOURCES and scope.}
\begin{enumerate}
\item Governing duty 4(7):
\path{<internal>/GOAL_alt.md}.
The H2 and static source entries of
\path{<internal>/STATE.tex}
are resolved at the finite prescribed-field and conditional
clock-surface scopes stated here.
\item Exact scalar law, trace and protected sectors:
\path{<internal>/rederivation_final3_bridge.tex},
Section 1, equations \texttt{br-polar}--\texttt{br-scalar-bound};
\path{<internal>/rederivation_final3_main_theorem.tex},
Section 4, both finite-collar proofs and the closed-parent conclusion.
These sources distinguish scalar identities, an ordered internal
lift and a physical constraint sector.
\item Positive ordinary-adjoint source and unchanged zero/one sectors:
\path{<internal>/DERIVATION.md},
Section 1, Theorem 1 and uniform packet dispersion. No inventory,
entropy or numerical claim from its other sections is used.
\item Gauss-row selection and auxiliary-algebra distinctions:
\path{<internal>/result.md},
Sections 1.1--1.5;
\path{<internal>/P16_derivation.md},
Sections 1 and 5. The graph flux, native commutator allocation,
rest-mode counterfamily and mean/strong selecting implications
are derived explicitly above, not imported as physical closure.
\end{enumerate}
\end{document}
